% Maze Lab - Project Proposal Report (Phase-I)
% Compile with pdflatex. 
% Ensure logo.png, member1.png, member2.png, and member3.png sit in the same directory.

\documentclass[12pt,a4paper]{article}

\usepackage[T1]{fontenc}
\usepackage[utf8]{inputenc}
\IfFileExists{newtxtext.sty}{\usepackage{newtxtext}}{\usepackage{mathptmx}}
\usepackage[a4paper,left=1.75cm,right=1.75cm,top=61.2pt,bottom=70pt,
            headheight=14pt,headsep=33.8pt,footskip=51.4pt]{geometry}
\usepackage{xcolor}
\usepackage{graphicx}
\usepackage{array}
\usepackage{tabularx}
\usepackage{fancyhdr}
\usepackage{eso-pic}
\usepackage{ragged2e}

% ---------- colours (sampled from template) ----------
\definecolor{gebanner}{RGB}{0,102,148}
\definecolor{gegold}{RGB}{255,201,38}
\definecolor{geheading}{RGB}{43,91,145}
\definecolor{gehint}{gray}{0.5}

% ---------- course & team definitions ----------
\newcommand{\courseinfo}{DSCPP-III-2026-T272}
\newcommand{\mone}{Tanmay Kumar}
\newcommand{\mtwo}{Sparsh Saxena}
\newcommand{\mthree}{Hemang Sethi}

% ---------- page style: header, page number, gold bar ----------
\pagestyle{fancy}
\fancyhf{}
\renewcommand{\headrulewidth}{0pt}
\fancyhead[L]{\fontsize{10.5}{12.6}\selectfont\courseinfo}
\fancyhead[R]{\fontsize{10.5}{12.6}\selectfont PROJECT PROPOSAL\hspace{2.6pt}}
\fancyfoot[R]{\fontsize{10}{12}\selectfont Page \thepage}
\AddToShipoutPictureBG*{%
  \AtPageLowerLeft{\color{gegold}\rule{\paperwidth}{4pt}}}

\setlength{\parindent}{0pt}
\setlength{\tabcolsep}{6pt}
\setlength{\topskip}{16.5pt}
\renewcommand{\arraystretch}{1}

% ---------- building blocks ----------
\newcommand{\banner}[3]{%
  \noindent\colorbox{gebanner}{\parbox[c][#3][c]{\dimexpr\textwidth-2\fboxsep\relax}{%
    \fontsize{#2}{#2}\selectfont\bfseries\color{gegold}#1}}\par}

\newcommand{\secthead}[2]{%
  {\fontsize{18}{21.9}\selectfont\color{geheading}#1\par}\vspace{1.6pt}%
  {\fontsize{10.5}{14.5}\selectfont\color{gehint}#2\par}%
  \vspace{\secgap}}
\newlength{\secgap}\setlength{\secgap}{5.7pt}
\newcommand{\afterbox}{\vspace{\aftergap}}
\newlength{\aftergap}\setlength{\aftergap}{16pt}

\newcolumntype{R}[1]{>{\RaggedRight\arraybackslash}p{#1}}
\newenvironment{answerbox}{%
  \noindent\begin{tabular}{|p{485.1pt}|}\hline
  \fontsize{11}{13.6}\selectfont\setlength{\parskip}{4pt}\rule{0pt}{22.5pt}\ignorespaces}{%
  \rule[-16pt]{0pt}{0pt}\\\hline\end{tabular}\par}

\newcommand{\memberrow}[4]{%
\hline
\fontsize{12}{14.5}\selectfont\rule{0pt}{22.5pt}\textbf{#1}\newline\textit{(Last Name, name: student ID: email, picture):}\rule[-65.1pt]{0pt}{0pt} &
\fontsize{12}{14.5}\selectfont\rule{0pt}{22.5pt}\textit{#3}\newline\textit{#4}\par\vspace{4.9pt}\noindent\raisebox{-4.35pt}{\includegraphics[width=1.5cm,height=1.5cm]{#2}}\\}

\begin{document}

% =====================  PAGE 1  =====================
\noindent\includegraphics[width=8cm]{logo.png}\par
\vspace{20.1pt}
\banner{PROJECT AND TEAM INFORMATION}{16.5pt}{22.5pt}
\vspace{22.8pt}

\setlength{\secgap}{11.4pt}
\secthead{Project Title}{(Try to choose a catchy title. Max 20 words).}
\begin{answerbox}
Maze Lab: Interactive Maze Generation and Comparative Algorithmic Pathfinding
\end{answerbox}
\setlength{\secgap}{5.7pt}
\vspace{23.8pt}

\noindent\begin{tabular}{|p{343.5pt}|R{129.2pt}|}
\hline
\fontsize{12}{14.5}\selectfont\rule{0pt}{22.5pt}\textit{Team Name:}\newline\textit{Team \#} &
\fontsize{12}{14.5}\selectfont\rule{0pt}{22.5pt}\textit{Team Blockbuster}\newline\textit{T272} \\
\memberrow{Team member 1 (Team Lead)}{member1.png}{Kumar, Tanmay -- 2510370439}{2510370439@geu.ac.in}
\memberrow{Team member 2}{member2.png}{Saxena, Sparsh -- 2510370313}{2510370313@geu.ac.in}
\memberrow{Team member 3}{member3.png}{Sethi, Hemang -- 2510370341}{2510370341@geu.ac.in}
\hline
\end{tabular}

\newpage
% =====================  PAGE 2  =====================
\banner{PROPOSAL DESCRIPTION (10 pts)}{17pt}{25.3pt}
\vspace{21.6pt}

\secthead{Motivation (1 pt)}{(Describe the problem you want to solve and why it is important. Max 300 words).}
\begin{answerbox}
Most core computer science courses teach data structures and algorithms through standalone coding tasks or static textbook diagrams[cite: 16]. Consequently, students often struggle to grasp how components like graph traversals, disjoint sets, or priority queues interact inside a single software application[cite: 16].

Maze Lab directly addresses this by offering a hands-on environment where players first navigate a generated maze manually and then evaluate their path against automated solvers[cite: 16]. Combining object-oriented system design with graph theory gives learners a direct, practical look at runtime performance, memory overhead, and heuristic search behavior[cite: 16].
\end{answerbox}
\afterbox

\secthead{State of the Art / Current solution (1 pt)}{(Describe how the problem is solved today (if it is). Max 200 words).}
\begin{answerbox}
Current learning tools usually split into simple text-only programs or web-based pathfinding visualizers[cite: 16]. Text exercises lack clear graphic feedback, while web visualizers rarely let users control a character through the maze[cite: 16]. Additionally, many existing codebases hard-code their algorithms, making it tricky to add or swap new solvers[cite: 16]. Maze Lab solves these issues by combining player controls with graph solvers using a clean, interface-driven C++ layout[cite: 16].
\end{answerbox}
\afterbox

\secthead{Project Goals and Milestones (2 pts)}{(Describe the project general goals. Include initial milestones as well any other milestones. Max 300 words).}
\begin{answerbox}
Our main goal is to build a flexible C++ application that generates valid mazes, supports manual player controls, and compares automated solving algorithms side by side[cite: 16].

Key milestones for implementation:
\par\vspace{2pt}
1. \textbf{Architecture Setup:} Define base \texttt{Cell} and \texttt{Maze} models alongside abstract \texttt{MazeGenerator} and \texttt{MazeSolver} interfaces[cite: 16].\newline
2. \textbf{Maze Generation:} Build an iterative Randomized DFS generator and a Kruskal's algorithm generator backed by Union-Find[cite: 16].\newline
3. \textbf{Player Mechanics \& Console UI:} Implement player navigation, wall collision detection, and ASCII terminal rendering[cite: 16].\newline
4. \textbf{Pathfinding Solvers:} Construct BFS and A* pathfinders that collect step counts, visited cell totals, and runtime metrics[cite: 16].\newline
5. \textbf{Validation \& Delivery:} Perform integration testing across all components, build a side-by-side benchmark mode, and complete final documentation[cite: 16].
\end{answerbox}

\newpage
% =====================  PAGE 3  =====================
\secthead{Project Approach (3 pts)}{(Describe how you plan to articulate and design a solution. Including platforms and technologies that you will use. Max 300 words).}
\begin{answerbox}
We are writing the program in C++ with an emphasis on modular object-oriented design[cite: 16]. The maze itself sits in a 2D grid of \texttt{Cell} structures, where removed walls correspond to graph edges[cite: 16].

To keep algorithms decoupled, we use abstract base classes under the Strategy Pattern[cite: 16]. \texttt{MazeGenerator} handles generation, while \texttt{MazeSolver} manages pathfinding[cite: 16]. Maze generation relies on explicit-stack DFS alongside Kruskal's algorithm built on a Disjoint Set Union class[cite: 16, 17]. Pathfinding tasks use BFS for shortest paths and A* guided by Manhattan distance heuristics[cite: 17]. Player logic lives separately in a \texttt{Player} class to verify movement before updating position[cite: 17]. Terminal output handles primary rendering, while SDL graphics remain an optional secondary goal[cite: 17].
\end{answerbox}
\afterbox

\secthead{System Architecture (High Level Diagram) (2 pts)}{(Provide an overview of the system, identifying its main components and interfaces in the form of a diagram using a tool of your choice).}
\begin{answerbox}
\begin{center}
\setlength{\unitlength}{0.9mm}
\begin{picture}(130,52)
  % Outer system boundary
  \put(0,0){\line(1,0){130}}
  \put(0,0){\line(0,1){52}}
  \put(130,0){\line(0,1){52}}
  \put(0,52){\line(1,0){130}}
  
  % Core Grid / Maze
  \put(45,34){\framebox(40,12){\shortstack{\textbf{Maze / Grid Module}\\\small(Cell, Grid, Boundaries)}}}
  
  % Strategy Interfaces
  \put(10,18){\framebox(34,10){\shortstack{\textbf{<<interface>>}\\\textbf{MazeGenerator}}}}
  \put(86,18){\framebox(34,10){\shortstack{\textbf{<<interface>>}\\\textbf{MazeSolver}}}}
  
  % Concrete Generators
  \put(5,2){\framebox(20,8){\small DFSGen}}
  \put(27,2){\framebox(20,8){\small KruskalGen}}
  
  % Concrete Solvers
  \put(83,2){\framebox(20,8){\small BFSSolver}}
  \put(105,2){\framebox(20,8){\small AStarSolver}}
  
  % Player & UI
  \put(52,18){\framebox(26,10){\shortstack{\textbf{Player / UI}\\\small(Terminal Render)}}}
  
  % Connectors
  \put(27,28){\line(2,1){18}}
  \put(103,28){\line(-2,1){18}}
  \put(15,12){\line(0,1){6}}
  \put(37,12){\line(-1,1){6}}
  \put(93,12){\line(0,1){6}}
  \put(115,12){\line(-1,1){6}}
  \put(65,28){\line(0,1){6}}
\end{picture}
\end{center}
\end{answerbox}
\afterbox

\secthead{Project Outcome / Deliverables (1 pts)}{(Describe what are the outcomes / deliverables of the project. Max 200 words).}
\begin{answerbox}
Key outcomes include a complete C++ terminal program, well-documented source code with clean commit records, automated unit tests, technical project documentation, and a live demonstration comparing human player paths against BFS and A* solutions[cite: 17].
\end{answerbox}
\afterbox

\secthead{Assumptions}{(Describe the assumptions (if any) you are making to solve the problem. Max 100 words)}
\begin{answerbox}
We assume the maze grid dimensions remain fixed during solving and that generated mazes are always simply connected with no loops[cite: 17]. We also assume the terminal host environment supports standard C++ I/O and ANSI color sequences[cite: 17].
\end{answerbox}
\afterbox

\secthead{References}{(Provide a list of resources or references you utilised for the completion of this deliverable. You may provide links).}
\begin{answerbox}
1. Cormen, T. H., et al. \textit{Introduction to Algorithms}. MIT Press[cite: 17].\newline
2. Gamma, E., et al. \textit{Design Patterns: Elements of Reusable Object-Oriented Software}. Addison-Wesley[cite: 17].\newline
3. ISO/IEC C++17 Standard Specifications[cite: 17].
\end{answerbox}

\end{document}